
% ============================================================
%  Seminar Paper Template — Topic 12
%  Numerical Simulation of Quantized Vortices with Small Core Size
%  Master-Seminar: Complex Quantum Systems and Vortex Dynamics
%  KIT — Fakultät für Mathematik, Sommersemester 2026
% ============================================================
\documentclass[12pt, a4paper]{article}

% ------------------------------------------------------------
%  Packages
% ------------------------------------------------------------
\usepackage[utf8]{inputenc}
\usepackage[T1]{fontenc}
\usepackage[english]{babel}

% Math
\usepackage{amsmath}
\usepackage{amssymb}
\usepackage{amsthm}
\usepackage{mathtools}

% Layout
\usepackage[
  top=2.5cm,
  bottom=2.5cm,
  left=2.8cm,
  right=2.8cm
]{geometry}
\usepackage{setspace}
\onehalfspacing

% Graphics & floats
\usepackage{graphicx}
\usepackage{float}
\usepackage{caption}
\usepackage{subcaption}

% References & links
\usepackage[
  backend=biber,
  style=alphabetic,
  sorting=nyt
]{biblatex}
\addbibresource{references.bib}   % <-- create this file with your .bib entries
\usepackage[
  colorlinks=true,
  linkcolor=black,
  citecolor=blue,
  urlcolor=blue
]{hyperref}

% Miscellaneous
\usepackage{microtype}
\usepackage{enumitem}
\usepackage{booktabs}

% ------------------------------------------------------------
%  Theorem environments
% ------------------------------------------------------------
\theoremstyle{plain}
\newtheorem{theorem}{Theorem}[section]
\newtheorem{lemma}[theorem]{Lemma}
\newtheorem{proposition}[theorem]{Proposition}
\newtheorem{corollary}[theorem]{Corollary}

\theoremstyle{definition}
\newtheorem{definition}[theorem]{Definition}
\newtheorem{example}[theorem]{Example}
\newtheorem{remark}[theorem]{Remark}

% ------------------------------------------------------------
%  Custom commands — add your own as needed
% ------------------------------------------------------------
\newcommand{\R}{\mathbb{R}}
\newcommand{\C}{\mathbb{C}}
\newcommand{\N}{\mathbb{N}}
%\newcommand{\i}{\text{i}}
\newcommand{\Eeps}{E_\varepsilon}
\newcommand{\psieps}{\psi_\varepsilon}
\newcommand{\ustar}{u^*}
\newcommand{\psistar}{\psi^*_\varepsilon}
\newcommand{\Wren}{W}
\newcommand{\Jac}{J\psi}
\newcommand{\sj}{j(\psi)}
\newcommand{\Lap}{\Delta}
\newcommand{\dd}{\,\mathrm{d}}
\newcommand{\norm}[1]{\left\lVert #1 \right\rVert}
\newcommand{\abs}[1]{\left\lvert #1 \right\rvert}
\newcommand{\inner}[2]{\left\langle #1,\, #2 \right\rangle}

% Shorthand for the symplectic matrix
\newcommand{\symp}{J_{\mathrm{symp}}}

% ------------------------------------------------------------
%  Title information
% ------------------------------------------------------------
\title{
  \Large\textbf{Numerical Simulation of Quantized Vortices\\
  with Small Core Size via Vortex Tracking}\\[0.5em]
  \large Master-Seminar: Complex Quantum Systems and Vortex Dynamics\\
  KIT --- Fakultät für Mathematik, Sommersemester 2026
}
\author{Thomas TODO-LASTNAME}   % <-- fill in your name
\date{\today}

% ============================================================
\begin{document}
% ============================================================

\maketitle
\thispagestyle{empty}

\section{Introduction}

Ginzburg-Landau Hamiltonian:
\begin{align}
    E_\varepsilon (u) = \int_\Omega \frac{1}{2} |\nabla u|^2 + \frac{1}{4 \varepsilon^2} (1- |u|^2)^2
\end{align}

When investigating this variational model and minimizes with the singular limit $\varepsilon \to 0$.

The corresponding Schrödinger flow $i \partial_t \psi = - \nabla_{L^2} E_\varepsilon (\psi)$
On smooth bounded domain with Neumann boundary conditions one arrives at:
\begin{align}
  i \partial_t \psi_\varepsilon = \Delta \psieps (x,t) + \frac{1}{\varepsilon^2} (1- |\psieps (x,t)|^2) \psieps(x,t)
\end{align}

Observables:
\begin{itemize}
  \item $|\psi|^2$
  \item Supercurrents $j(\psi) = \frac{1}{2i}(\overline{psi} \nabla - \psi \nabla \overline{\psi})$
  \item Vorticity: $J \psi = \frac{1}{2} \nabla \times j(\psi)$
\end{itemize}

For vortices:
$J \psieps (t) \to \pi \sum^N_{j=1} d_j \delta_{a_j (t)}$
The singular points solve:
\begin{align}
  \dot c_j (t) = - \frac{1}{\pi} d_j J \nabla_{a_j} W(a(t), d) \\
  a_j(0) = a^0_j
\end{align}


We introduce the harmonic map $u^* (a,d) \in W^{1,p}(\Omega , \mathbb S^{-1})$.
What is it?
Given vortex positions $a$ and degrees $d$ we want a map $u^*: \Omega \to \C$ that:
\begin{itemize}
  \item lives on the unit circle away from vortex locations
  \item has exactly the right topological winding number $d_j$ around each $a_j$
  \item satisfies homogeneous Neumann boundary conditions $\partial_\nu u^* = 0$ on $\partial \Omega$
\end{itemize}

The following conditions uniquely determine the supercurrents $j(u^*)$ which in turn determine $u^*$
\begin{align}
  \nabla j(u^*) =0\\
  \nabla \times j(^*) = 2\pi \sum^N_{j=1} d_j \delta_{a_j} \\
  \nu j(u^*) =0 \text{ on } \partial \Omega
\end{align}
Where $\nu$ is the direction of the outward normal of $\partial \Omega$
$u^*$ is then given by
\begin{align}
  u^* (x; a,d) = \exp{i H(x)} \prod^N_{j=1} \frac{x-a_j}{|x-a_j|}^{d_j}
\end{align}

Where $H:\R^2 \to \R^2$ is some harmonic function.

Boundary terms in the renormalized energy $R(x; a,d)$ which solves
\begin{align}
  \Delta R = 0 \text{, in \Omega} \\
  R = - \sum^N_{j=1} d_j \log |x-a_j| \text{, on } \partial \Omega
\end{align}

The supercurrents of $u^*$ satisfies:
\begin{align}
  j(u^*) = - \nabla \times G, \quad G(x) = \sum_j d_j \log |x-a_j| + R(x; a,d) 
\end{align}
This equation is important to compute the gradient of $W$
$u^*$ has singularities at the vortex locations $a_j$. This is the reason we need to smooth it out before using it as an initial condition for the PDE.



The renormalized energy:
Why do we need it?
The Ginzburg-Landau energy diverges as $\varepsilon \to 0$. This divergence comes purely from the self-energy of each vortex core and carries no information about vortex interactions.
The renormalized energy is what remains after subtracting the divergence:
\begin{align}
  W(a,d) =  - \pi \sum_{1 \leq i \neq j \leq N} d_i d_j \log |a_i - a_j - \pi \sum_j d_j R(a_j ; a, d)
\end{align}
The first terms captures vortex-vortex interaction: same sign vortices repel, oppositve sign attract
The second sum captures vortex-boundary interaction through R, the harmonic correction that enforces the Neumann boundary condition

\begin{definition}
  A familiy of iniital conditions $(\psieps^0)_{\varepsilon>0}$ is siad to be well-prepared if it satisfies the following assumptions,
  for some constant $C>=; 0<\alpha<1$ and $\varepsilon$ small enough:
  \begin{enumerate}
    \item There exists $N$ vortices with positions
    $a^0 = (a^0_j)_{j=1, \ldots , N} \in \Omega^{*N}$
    and degrees $d= (d_j)_{j=1, \ldots, N} \in \{\pm 1\}^N$
    such that
    \begin{align}
      ||J\psieps^0 - \pi \sum^N_{j=1} d_j \delta_{a_j^0}|| \leq C \varepsilon^\alpha 
    \end{align}
    and the vortices are distant enough
    \item the energy of $\psieps^0$ is close to be optimal:
    \begin{align}
      E_\varepsilon (\psieps^0) \leq W_\varepsilon (a_0, d) + C \varepsilon^{1/2}
    \end{align}
  \end{enumerate}
\end{definition}


\section{Method}
\subsection{Motivation}
Why not just solve the PDE directly?
\begin{itemize}
  \item The standard approach: time-splittinng + Fourier /FEM in space
  \item The core problem: resolving the vortex of size $\varepsilon$ requires mesh size $h<\varepsilon$
  \item Adaptive FEM helps, but still expensive, the method in this paper shifts the bottleneck entirely
\end{itemize}



\subsection{Algorithm}

Step 1: Solve $1D$ ODE for $f_varepsilon$ compute initial phase $H$
Step 2: Evolve vortex positions via Hamiltonian ODE (RK4, $\delta t = 10^{-3}$)
Step 3: The radial prodfile $f_\varepsilon$ solved once, scales with $r_0/\varepsilon$

Key insight: $\epsilon$ only enters through the 1D preprocessing - it is no longer a discretization bottleneck

\subsection

\section{Error Analysis}

\begin{lemma}
  Let $\Omega \subset \R^2$ be the unit disk, let $\{a(t)\}_{t\geq 0}$ denote the exact solution to the Hamiltonian mechanics.
  Let $b(t)$ be the approximated ODE trajectory with the same degrees and initial positions $b^0$ with harmonic polynomials of degree up to $n$
  and time step $\delta_t$. Suppose that for some $T>0$,
  \begin{align}
  \rho_T = \min_{t\leq T} \min_{j\neq k} \{ |b_j (t) - b_k  (t)|, \text{dist} (b_k(t) , \partial \Omega), |a_j (t) - a_k  (t)|, \text{dist} (a_j(t) , \partial \Omega)\}>0
  \end{align}
  is such that $0 < \rho_T < 1$. Then we have:
  \begin{align}
    |a(t) - b(t)| \leq C (|a^0 -b^0| + (1- \rho_T)^n n^{1-l}) + C \delta t^4
  \end{align}
\end{lemma}

In order to prove this we need 2 other lemmas 

\begin{lemma}
  Let $b^0, a^0 \in \Omega^{*N}$, for fixed $d \in \{\pm1\}^N$ and let $\rho(a) := \min_{j\neq l} \{\text{dist}(a_j, \partial\Omega) , |a_j - a_l|\}$.
  Suppose that $\rho_T:= \text{inf}_{t\in [0,T]} \leq 0$
  \begin{align}
  |a(t) - b(t)| \leq  C (|a^0-b^0| + \frac{\text{sup}_{s\in [0,T]}||\mathbb{P}^\vert_n g_{a(s)}||_{L^2(\partial \Omega)}}{\rho^{3/2}_T})\exp{C t /\rho^{5/2}_T}
  \end{align}
  Where $g_a(x) = - \sum^N_j d_j \log |x-a_j|$
\end{lemma} 

The time integration errors follows from the previous bounds and standard considerations on the numerical integration of ODEs

\begin{lemma}
  Let $b$ be the exact solution to $\dot{b} = F^n (b)$ and let $(b^k_{\delta t})_{k \delta t \leq T}$ be its numerical approximation with RK4 method and time step $\delta t$
  \begin{align}
    b^0_{\delta t} = b^0\\
    b^0_{\delta t} = b^{k-1}_{\delta t} + \delta t \Phi_{\delta t , n} (b^{k-1}_{\delta t})
  \end{align}
Where $\Phi_{\delta t, n}$ is the RK4 increment function associated with the approximate forcing term $F^n$.
Assume that:
\begin{align}
  0 < \rho_T = \min (\inf_{0 \leq t \ leq T} \rho(b(t)), \min_{k \delta t \leq T} \rho (b^k_{\delta t}))
\end{align}
Then there exists constants $C_{\rho_T}>0$ and $\Lambda_{\rho_T}$, that depends on $\rho_T$ but not on $n, b$ or $\delta t$ such that
the following in time error holds:
\begin{align}
  \max_{k \delta t < T} |b(k \delta t) - b^k_{delta_t}| \leq \delta t^4 \frac{C_{\rho_T}}{\Lambda_{\rho_T}} (\exp (\Lambda_{\rho_T} T)- 1)
\end{align}
\end{lemma}

% ============================================================
%  Bibliography
% ============================================================
\newpage
\printbibliography